\subsection{HD-I09: nonfaithful coordinate powers and invisible branch points}
\label{hd-i09}

Replace the coordinate-line character $\chi$ by $\chi^k$, and put
$g_0=\gcd(d,k)$, $d'=d/g_0$, $e'=\gcd(e,d')$. There is an exact
surjective group map
\[
 G(d,e,m)\longrightarrow G(d',e',m),\qquad
 \operatorname{diag}(\zeta_d^{a_i})P_\pi\longmapsto
 \operatorname{diag}(\zeta_{d'}^{(k/g_0)a_i})P_\pi.
\]
Multiplication by $k/g_0$ is a unit modulo $d'$. An exponent vector
modulo $d'$ lifts to one with total zero modulo $e$ exactly when its
total is zero modulo $\gcd(e,d')$: changing one lift by multiples of
$d'$ solves precisely this congruence. This proves surjectivity.

For $N_k=\ker\chi^k$ the exact inducing map again gives the cyclic
tower $X_k=Z/N_k\to Y=Z/H$ of degree $d'$. The primitive factor of
$J(X_k)$ maps to the power-character multiplicity factor by the actual
pullback isogeny with polarization multiplier $|N_k|$, using the same
self-adjoint primitive projector as HD-I02. In particular, this is an
exact period factorization for the nonfaithful representation.

An original colored cycle $(\ell,s)$ has new color
\[
 s'=(k/g_0)s\pmod{d'},\qquad 0\le s'<d'.
\]
It contributes to the branch divisor of $X_k\to Y$ precisely when
$s'\ne0$. If $r'_a$ counts these new nonzero colors, the signature
conditions are
\[
 g(Y)=1-m+\frac12\sum(\ell-1),\quad
 r'=\sum_a r'_a=2(n-g(Y)+1),\quad
 \sum_a(2a-d')r'_a=0,
\]
together with the original product-one and full-generation Nielsen
conditions. The primitive period rank is at most $3n$, or $3n-1$
when $e'<d'$, by the same quotient and codifferential morphisms.

For $d=3,4,6$ the additional imaginary quadratic case is
$d=6$, $k=2$ or $4$, for which $d'=3$. Original color $3$ is then
invisible. Powers of order at most two have character field
$\mathbb Q$ and do not themselves provide this quadratic Weil action.
For the cubic power case, the sharp maxima at $n=5$, for every fixed
$m\ge9$, are
\[
\begin{array}{c|rrrr}
 \text{original }e&1&2&3&6\\\hline
 \text{image }e'&1&1&3&3\\
 \max\dim\operatorname{im}(\mathrm{period})&14&14&15&15.
\end{array}
\]

For attainment, begin with the rank-fifteen unramified cubic genus-six
covers or the rank-fourteen cubic genus-five covers branched with
exponents $1,2$ at $P,Q$. Take the normalized fibre product with a
connected double cover of the same base. Coprime degrees make the
result a connected cyclic six-cover. For $e=6$ use an \'etale double
cover; for $e=3$ use one branched at two new points. The six-cover
exponents are respectively absent or $3,3$. For $e=2$ use an \'etale
double cover, giving exponents $4,2$ at $P,Q$; for $e=1$ use one
branched at $P,Q$, giving exponents $1,5$. Such double covers exist
because a square root of $\mathcal O(P+Q)$ exists in
$\operatorname{Pic}^1$ over $\mathbb C$, and positive-genus curves
have nonzero two-torsion.

These are dominant finite root choices over the cubic-cover moduli,
with two extra free markings for $e=3$. A general degree-$m$ pencil
as in HD-I08 gives full monodromy $G(6,e,m)$: the connected
colored-transposition graph and its diagonal rotations by $e$ are
exactly the generators in that proof. Forgetting the auxiliary data
is dominant, so the certified ranks persist. Thus invisible branch
points may enlarge the original Hurwitz source while leaving its
primitive period image governed by the displayed cyclic quotient.

\subsection{HD-I10: sharp maxima for every every requested dimension}
\label{hd-i10}

For every $n\ge5$ and
$m\ge\lceil(n^2-2n+3)/2\rceil$, the maximum period-image dimension
over branch components of the standard $G(d,e,m)$ representation is
exactly $3n$ for $e=d$, and $3n-1$ for $e<d$. We prove attainment by
an exact one-node induction from the certificates in HD-I06--I07.

\paragraph{The three new polar directions.}
Add one node $(u,v)$ to the rational normalization of an old nodal
model of genus $g$, keeping the branch divisor $R$ empty or equal to
$P+Q$. Old eigensections embed in the new spaces: they have zero
residues at the new node. The old quadratic target
$H^0(\omega_{\rm old}^2(R))$ has codimension three in
$H^0(\omega_{\rm new}^2(R))$. The quotient consists exactly of the
common new double-pole coefficient and the two simple-pole
coefficients. Matching the double coefficients is the new
dualizing-square gluing condition; vanishing of all three means
regularity at $u,v$, hence membership in the old target.

Let the new forward gluing be $\lambda$. Choose new forward and inverse
sections with residues $(1,-\lambda)$ and $(1,-\lambda^{-1})$.
Their product has common double coefficient one. Products of the new
forward section with old inverse sections $g_i$ have simple
coefficient vectors
\[
 (g_i(u),-\lambda g_i(v)),
\]
in the fixed normalization frames. Two linearly independent old
sections have nonconstant rational ratio, so these two vectors are
independent for a general pair $(u,v)$. Therefore the new products
span all three quotient directions. Surjectivity of the old product
map implies surjectivity of the new one, increasing rank by three.

\paragraph{Unramified covers.}
Keep all node gluings equal to $\lambda=\zeta_d$. With an old node
$(a_*,b_*)$, choose
\[
 f_{\rm new}=
 \left(\frac1{z-u}-\frac\lambda{z-v}
       -\frac1{z-a_*}+\frac\lambda{z-b_*}\right)dz,
\]
and the same expression with $\lambda^{-1}$ for the inverse section.
These are regular at infinity and have the required residues. An old
node keeps the line bundle of exact order $d$ and its cover connected.
Starting with genus six and rank fifteen gives genus $n+1$ and rank
$3n$ for every $n\ge5$. The torsor smoothing, constant dimensions and
base-change argument of HD-I06 give smooth covers of the same rank.

\paragraph{Two branch points.}
Keep $P=0$, $Q=\infty$, $c=7^d$, $\lambda=\zeta_d7^e$ and the root
section $z^e$ of HD-I07. Choose the new pair as $(u,v)=(u,cu)$.
Then $v^e=\lambda^du^e$, preserving the exact root relation.
The new sections are
\[
 f_{\rm new}=\left(\frac1{z-u}-\frac\lambda{z-v}\right)dz,
 \quad
 g_{\rm new}=\left(\frac1{z-u}-\frac{\lambda^{-1}}{z-v}
                     -\frac{1-\lambda^{-1}}z\right)dz.
\]
The first has only its allowed pole at infinity; the second has the
allowed pole at zero and is regular at infinity. The two simple-pole
vectors remain independent for some $u$. Otherwise the nonconstant
ratio $R$ of two old inverse sections would satisfy $R(cz)=R(z)$.
Since $c$ has infinite multiplicative order, this would give an
infinite fibre of a nonconstant rational function, a contradiction.
Choose $u$ away from the finitely many determinant zeros and node or
branch coincidences. Starting with genus five and rank fourteen
therefore gives genus $n$ and rank $3n-1$. The connected marked-root
smoothing in HD-I07 gives smooth covers attaining this rank.

\paragraph{The requested degree and full monodromy.}
The pencil argument of HD-I08 works in every genus for $m\ge g+3$:
the bad-line-bundle locus has dimension at most $2g-m+2\le g-1$.
Moreover
\[
 \left\lceil\frac{n^2-2n+3}{2}\right\rceil\ge n+4
 \quad(n\ge5),\qquad
 n^2-2n+3-2(n+4)=(n-5)(n+1).
\]
Thus the required $m$ supplies simply branched pencils for both
genera. The connected colored-transposition graph and diagonal
rotation argument gives full $G(d,e,m)$, exactly as in HD-I08.
Forgetting the pencil is dominant, so the cyclic ranks are attained
by full Hurwitz components. The upper bounds of HD-I04 prove the
stated maxima. Since $3n<n^2$, none dominates.

The coprime-cover lifts of HD-I09 also extend verbatim: for the cubic
coordinate powers of $G(6,e,m)$ the sharp maxima are $3n$ for $e=3,6$
and $3n-1$ for $e=1,2$, throughout the same degree range.
